\section{Algorithm for best Egyptian decomposition}
\phantomsection
\label[appendix]{app:best-egyptian-decomposition}

Fibonacci's greedy method shows that every irreducible proper fraction $\dfrac{b}{a}$ can be decomposed into at most $b$ Egyptian fractions. The \enquote{best} decomposition should contain the fewest Egyptian fractions possible. If several decompositions have the same number of terms, the one with the smaller denominator is better. We maintain a set $B = \{\dfrac{1}{q_1}, \dfrac{1}{q_2}, \ldots, \dfrac{1}{q_m}\}$ to record the best candidate and initialize $B=\varnothing$. The following are the steps of the algorithm:

\begin{enumerate}[1)]
  \item Decompose $\dfrac{b}{a}$ into no more than $n=b$ Egyptian fractions. First, try $\dfrac{b}{a}=\dfrac{1}{q}+x$, where $q=\lfloor\dfrac{a}{b}\rfloor+1$, one more than the quotient of $a$ divided by $b$. This step is the same as Fibonacci's method.
  \item Simplify the problem to decomposing $x$ into no more than $n-1$ Egyptian fractions. Add $\dfrac{1}{q}$ to $Q$, so $Q=\{\dfrac{1}{q}\}$. Then calculate $x=\dfrac{b}{a}-\dfrac{1}{q}=\dfrac{bq-a}{aq}$ and reduce the result to lowest terms as $\dfrac{b'}{a'}$ (use the Euclidean algorithm to find the greatest common divisor of the numerator and denominator, then divide both by it).
  \item If $b'=1$ and $a'>q$, we have found a \emph{candidate} decomposition. We restrict $a'>q$ because any permutation of the unit fractions in $Q$ gives the same solution. To avoid duplicates, we require $\dfrac{1}{q_0}>\dfrac{1}{q_1}>\cdots>\dfrac{1}{q_m}$. When a new solution is found, add $\dfrac{1}{a'}$ to $Q$, then compare $Q$ with $B$. If $B=\varnothing$, $|Q|<|B|$, or the maximum denominator in $Q$ is smaller than that in $B$, replace $B$ with $Q$ and continue searching.
  \item If $b' \ne 1$, we append $\dfrac{b'}{a'}$ together with $Q$ as a pair to the \emph{tail} of a list\footnote{Actually a queue in computer programming} $C$, such that we will decompose $\dfrac{b'}{a'}$ on top of $Q$ in the future.
  \item Increase $q$ to $q'=q+1$, and attempt to decompose $\dfrac{b}{a}=\dfrac{1}{q'}+x$. This decomposition is less greedy than Fibonacci's choice of $q$, so the remainder $x$ is larger. Check whether $n\dfrac{1}{q'}<\dfrac{b}{a}$; if it is, $n$ fractions no larger than $\dfrac{1}{q'}$ cannot reach $\dfrac{b}{a}$, so decomposing $x$ further is impossible. Stop and remove a pair from the \emph{head} of $C$. This pair contains $\dfrac{b'}{a'}$ and $Q$, from which the search continues. Otherwise, if $n\dfrac{1}{q'}\geq\dfrac{b}{a}$, return to step 2 and continue attempting to decompose $x$ into no more than $n-1$ Egyptian fractions.
  \item We try every candidate from $q, q' = q + 1, q'' = q + 2 \cdots $ to $q + i \geq \dfrac{an}{b}$. Because the Egyptian fractions in $Q$ are descending (their denominators are ascending), the $x$ in $\dfrac{b}{a} = \dfrac{1}{q+i} + x$ cannot be decomposed into $n-1$ Egyptian fractions being less than $\dfrac{1}{q+i}$.
  \item Each time a pair $\dfrac{b'}{a'}$ and $Q$ is removed from the head of $C$, and $Q$ is not empty, we (a) compare $q_0$, the denominator of the smallest fraction in $Q$, with $q=\lfloor\dfrac{a'}{b'}\rfloor+1$, and use the larger value for further decomposition; (b) compare $Q$ with $B$. If $B$ is not empty and contains fewer Egyptian fractions, $Q$ cannot be optimal. Because candidates containing more fractions are appended to the tail of $C$, all remaining candidates are no better, so $B$ is optimal and the algorithm terminates.
\end{enumerate}

The following is the formal description of the algorithm. The reader who knows computer programming may implement it.

\begin{algorithmic}[1]
\Function{Decompose}{$\dfrac{b}{a}$}
  \If{$b = 1$}
    \State \Return $\{\dfrac{1}{a}\}$
  \EndIf
  \State $n = b, B = \varnothing$
  \State $C = \{(\dfrac{b}{a}, \varnothing)\}$
  \While{$C \ne \varnothing$}
    \State $(\dfrac{b}{a}, Q) \gets $ \Call{Pop}{$C$}  \Comment{poll from the head of $C$}
    \State $n' = n - |Q|$                            \Comment{decompose at most $n'$ Egyptian fractions}
    \State $q = \lfloor \dfrac{a}{b} \rfloor + 1$
    \If {$\varnothing \ne Q = \{\dfrac{1}{q_0}, \cdots \}$ and $q_0 \geq q$}  \Comment{The first fraction in $Q$ is less.}
      \State $q = q_0 + 1$
    \EndIf
    \While{$q \leq \dfrac{an'}{b}$}
      \If{$B \ne \varnothing$ and $|Q| \geq |B| - 1$} \Comment{Candidates in $C$ are worse than $B$}
        \State \Return $B$
      \EndIf
      \State $\dfrac{b'}{a'} = $ \Call{Reduce}{$\dfrac{b}{a} - \dfrac{1}{q}$} \Comment{reduce to lowest terms}
      \If{$b' = 1$ and $a' > q$}
        \State compare $\{\dfrac{1}{a'}, \dfrac{1}{q}\} \cup Q$ and $B$, record the better as $B$
      \Else
        \State \Call{Push}{$(\dfrac{b'}{a'}, \{\dfrac{1}{q}\} \cup Q), C$} \Comment{Append to the tail of $C$}
      \EndIf
      \State $q \gets q + 1$
    \EndWhile
  \EndWhile
  \State \Return $B$
\EndFunction
\end{algorithmic}
